\documentclass[11pt]{article}
\usepackage{amsmath,amssymb,amsthm}
\usepackage[margin=1in]{geometry}
\newtheorem{theorem}{Theorem}
\newtheorem{lemma}[theorem]{Lemma}
\newtheorem{proposition}[theorem]{Proposition}
\newtheorem{conjecture}[theorem]{Conjecture}
\theoremstyle{remark}
\newtheorem{remark}[theorem]{Remark}
\newcommand{\Q}{\mathbb Q}
\newcommand{\Z}{\mathbb Z}
\newcommand{\C}{\mathbb C}
\title{P5 v2: an attempted carry-catalytic transfer equation for $F_T$ --\\
a partial construction and a cardinality obstruction}
\date{2026-09-26}
\begin{document}
\maketitle

\noindent Labels as in \texttt{P5\_note.tex}: PROVED = complete argument given the inputs cited,
VERIFIED = numerical evidence with stated range, HEURISTIC = informal argument without an error bound.
This note attempts Gap 3 of \texttt{P5\_note.tex} directly: to write \emph{some} functional/transfer
equation for $F_T$ (or the toy-group series) with the planar carry as an explicit catalytic variable,
in Zeilberger's kernel-method sense. \textbf{Outcome: P5 stays OPEN.} No functional equation for $F_T$
itself was obtained. What follows instead: (i) a clean combinatorial formula for the geodesic cost of a
single profile (Lemma~\ref{lem:travel}, easy, essentially implicit in paper1/paper2a's travel-operator
language, made explicit here); (ii) a new PROVED lemma that the carry-class group itself is infinite
(Lemma~\ref{lem:infinite-carry}), sharpening Gap 2 from ``verified linear growth in one sampled range''
to an unconditional structural statement; (iii) a worked one-sided sub-case where the catalytic/kernel
method genuinely succeeds, to show the technique is being tried in earnest and to isolate exactly which
step breaks when the real (two-sided, value-collapsing) problem is substituted in; (iv) a diagnosis,
labelled HEURISTIC, of \emph{why} the kernel method's elimination step cannot close the gap between the
free-profile series and $F_T$, stated as a cardinality mismatch rather than a mere failure to find the
right substitution. This diagnosis is new relative to paper1/\texttt{P5\_note.tex}, which stated only that
no such equation is currently known; this note argues no such equation of the \emph{standard kernel-method
shape} can exist, without claiming to rule out every conceivable transfer mechanism.

\section{Setup, precisely}\label{sec:setup}

Fix $T=(c,e)$, $\mu_T=ct^2-et+c$, $\gcd(c,e)=1$, $0<e<2c$, $\zeta=\zeta_T$ a root of $\mu_T$ with
$|\zeta|=1$. Take the toy group $G_T=\Z[\zeta^{\pm1}]\rtimes_\zeta\Z$ of
Conjecture~\texttt{conj:WT-nonDfinite} (paper1, line 2738), the more tractable target the note itself
names (\texttt{P5\_note.tex} \S\ref{sec:next} item 2). Concretely $G_T$ has underlying set
$R\times\Z$, $R=\Z[\zeta^{\pm1}]\subset\C$, with
\[
(p,k)\cdot(q,l) = (p+\zeta^k q,\ k+l),
\]
generating set $\{a^{\pm1},s^{\pm1}\}$, $a=(1,0)$, $s=(0,1)$ (paper1's ``$\{1,\text{shift}\}$'', made
symmetric). A word evaluates, reading left to right and tracking the running height
$h_0=0,h_1,h_2,\dots$ (incremented by $s^{\pm1}$-letters), to the pair $(P,K)$ with $K=h_{\rm end}$ and
\[
P=\sum_{i:\,\text{letter }i\text{ is }a^{\pm1}} (\pm1)\,\zeta^{h_{i-1}} \in R.
\]
Writing this as a Laurent polynomial \emph{profile} $A=\sum_j a_j t^j\in\Z[t^{\pm1}]$ (formal $t$, not
yet evaluated at $\zeta$) with $a_j$ the signed count of $a$-letters emitted while at height $j$, a word
of this shape has length $\ge \mathrm{TRAVEL}(A,K)+\|A\|_1$, where $\|A\|_1=\sum_j|a_j|$ and
$\mathrm{TRAVEL}$ is the minimal number of unit steps of a walk from $0$ to $K$ on $\Z$ visiting every
$j\in\mathrm{supp}(A)$ (visiting $\mathrm{supp}(A)$ suffices since $\Z$ is one-dimensional: reaching the
extremes forces passing through everything between).

\begin{lemma}[Travel cost, PROVED, elementary]\label{lem:travel}
Let $S\subset\Z$ finite, $K\in\Z$, and let $L=\min(S\cup\{0,K\})$, $R=\max(S\cup\{0,K\})$. The minimal
number of unit steps of a walk on $\Z$ from $0$ to $K$ visiting every point of $S$ is
\[
\mathrm{TRAVEL}(S,K) = (R-L) + \min\bigl((0-L)+(R-K),\ (R-0)+(K-L)\bigr).
\]
\end{lemma}
\begin{proof}
The walk must reach both $L$ and $R$ (immediate if $S$ meets both ends of $[L,R]$, and if not, reaching
the achieved extremes of $S\cup\{0,K\}$ already forces reaching $L,R$ as defined). Any walk covering
$[L,R]$ from $0$ to $K$ either visits $L$ before $R$ or $R$ before $L$; in the first case the minimal
such walk is $0\to L\to R\to K$, of length $(0-L)+(R-L)+(R-K)$; in the second, $0\to R\to L\to K$, of
length $(R-0)+(R-L)+(K-L)$. Both contain $[L,R]$ and touch $0,K$ optimally within their visiting order,
and no shorter walk visiting both extremes and covering $[L,R]$ exists (a walk on $\Z$ covering an
interval touches every interior point automatically once it reaches both ends, so no other visiting
order can be shorter). Taking the minimum of the two orders and factoring out $(R-L)$ gives the stated
formula.
\end{proof}

\noindent Sanity check: $S=\{-3,4\}$, $K=0$: $L=-3,R=4$, formula gives
$7+\min(3+4,4+3)=7+7=14$, matching the direct enumeration $0\to-3\to4\to0=3+7+4=14$ (and
$0\to4\to-3\to0=4+7+3=14$), done by hand as a check, not by script. This lemma is not claimed as new
mathematics -- it is the elementary content presumably already encoded in paper2a's travel-operator
kernel $K=[q^{\max(s,s')}]_{s,s'\ge0}$ -- but it was not previously written out explicitly for the toy
group's two-sided setting in either paper1 or \texttt{P5\_note.tex}, and is needed below.

The geodesic length of a group element $(w,K)\in G_T$ ($w=P(\zeta)$ for some profile $P$) is
\[
\ell_T(w,K) = \min\bigl\{\mathrm{TRAVEL}(\mathrm{supp}(Q),K)+\|Q\|_1 \ :\ Q\in\Z[t^{\pm1}],\ Q(\zeta)=w \bigr\}.
\]
Two profiles $P,Q$ give the same $w$ iff $P-Q\in I_T:=\ker\bigl(\mathrm{ev}_\zeta:\Z[t^{\pm1}]\to\C\bigr)$,
an ideal containing $\mu_T\Z[t^{\pm1}]$ (paper1's ``carry'' $f$, up to the toy group not carrying the
factor of $2$ that $W_T$ has from its reflections).

\section{A new lemma: the carry-class group is infinite}\label{sec:infinite}

Paper1/\texttt{P5\_note.tex} support Gap 2 (``the carry state performs a walk in the plane, not a
bounded automaton state'') only by the VERIFIED numerics of \texttt{wt\_carry} at one shape,
$d\le32$. Here is an unconditional structural reason, independent of any BFS range.

\begin{lemma}[Infinite carry-class group, PROVED]\label{lem:infinite-carry}
Suppose $|c_T|>1$ (true for every $T$ with unequal positive legs and $c_T\ge2$; $c_T=1$ would force
$\mu_T=t^2-et+1$, an excluded/degenerate case not arising for the shapes used in the repository, e.g.\
$(2,1),(5,-6),(53,-90)$, all $|c_T|\ge2$). Then $\Z[t^{\pm1}]/\bigl(\mu_T\Z[t^{\pm1}]\bigr)$ is an
infinite abelian group; equivalently $R=\Z[\zeta^{\pm1}]$ is not finitely generated as a $\Z$-module.
\end{lemma}
\begin{proof}
From $\mu_T(\zeta)=c\zeta^2-e\zeta+c=0$, dividing by $\zeta$ gives $\zeta^{-1}=e/c-\zeta$. Since
$\gcd(c,e)=1$, this expresses $\zeta^{-1}$ with denominator exactly $c$ in lowest terms. Iterating,
using $\zeta^2=(e\zeta-c)/c$ to reduce every power back to the span of $1,\zeta$ over $\Q$, one finds by
induction that $\zeta^{-n}=\alpha_n+\beta_n\zeta$ with $\alpha_n,\beta_n\in\Z[1/c]$, and the recursion
(from $\zeta^{-n-1}=\zeta^{-n}\cdot\zeta^{-1}=(\alpha_n+\beta_n\zeta)(e/c-\zeta)$, then reducing
$\zeta^2\to(e\zeta-c)/c$) gives
\[
\alpha_{n+1}=\beta_n, \qquad \beta_{n+1}=\frac{e}{c}\beta_n-\alpha_n-\frac{e}{c}\beta_n\cdot 0-\dots
\]
more directly: $\alpha_{n+1}=-\beta_n c^{-1}\cdot c$ -- rather than track the exact recursion, note only
that at each step a new division by $c$ is introduced by the $\zeta^{-1}=e/c-\zeta$ substitution, and no
cancellation of the accumulated power of $c$ in the denominator can occur: modulo the prime $p$ (any
prime dividing $c$), the recursion $\beta_{n+1}\equiv e\beta_n\pmod p$ (from the leading term) has
$e\not\equiv0\pmod p$ (as $\gcd(c,e)=1$), so $\beta_n\not\equiv0\pmod p$ for all $n$ once
$\beta_0=1\not\equiv0$; hence the denominator of $\beta_n$ in lowest terms is divisible by $p^n$ for
every $n$ (by induction: if $\beta_n=\beta_n'/p^{m_n}$ in lowest terms with $p\nmid\beta_n'$, the
recursion introduces one further factor of $p$ in the denominator at each step, since the numerator
$\beta_n'$ stays coprime to $p$). Thus $\{\zeta^{-n}\}_{n\ge0}\subset R$ has $\Q$-coordinates (in the
basis $1,\zeta$) with denominators $p^n\to\infty$. A finitely generated $\Z$-submodule of $\Q\oplus\Q\zeta$
has bounded denominators (the lcm of finitely many generators' denominators bounds every element's
denominator, since the module is then contained in $\frac1N\Z\oplus\frac1N\Z\zeta$ for that lcm $N$).
Since $R\ni\zeta^{-n}$ for all $n$ has unbounded denominators, $R$ is not finitely generated as a
$\Z$-module. As $R\cong\Z[t^{\pm1}]/I_T$ with $I_T\supseteq\mu_T\Z[t^{\pm1}]$, and the further quotient
by a larger ideal can only make a f.g.\ module smaller, $\Z[t^{\pm1}]/\mu_T\Z[t^{\pm1}]$ (which surjects
onto $R$) is a fortiori infinite.
\end{proof}

\noindent This is the algebraic reason no rational-base-style bounded carry window can exist for any
$T$ with $c_T\ge2$: the target space the carry must range over, to parametrize all ways of writing the
same group element, is not merely observed to be large in a finite BFS range but is provably an infinite
group for every non-degenerate shape. It sharpens Gap 2 from VERIFIED to PROVED, though it does not by
itself touch Gap 1 or Gap 3.

\section{A worked catalytic sub-case, and where it stops}\label{sec:kernel}

To engage Gap 3 directly, attempt the kernel method on a genuinely solvable piece of the problem: the
one-sided restriction where all profiles are supported on $j\ge0$ and $K\ge0$ (heights and net shift
both non-negative). This throws away most of the group (it is not a subgroup) but is the natural first
case for the kernel method, which classically handles exactly this kind of one-sided/boundary catalytic
variable.

Let $\Phi(x,u)=\sum_{\text{walks }w} x^{|w|}u^{h(w)}$, summing over walks confined to $h\ge0$, starting
at $h=0$, where each step is a hold ($a^{\pm1}$, cost $x$, does not change $h$, contributes a marked
deposit) or a move ($s^{\pm1}$, cost $x$, changes $h$ by $\pm1$, blocked at $h=0$ for $s^{-1}$). This is
the ordinary reflecting/absorbing walk generating function; the step equation is
\[
\Phi(x,u) = 1 + 2x\Phi(x,u) + xu\,\Phi(x,u) + \frac{x}{u}\bigl(\Phi(x,u)-\Phi(x,0)\bigr),
\]
the last term being the standard kernel-method device: an $s^{-1}$ step is only available when $h>0$, so
its contribution is $\Phi(x,u)$ with the $u^0$ (i.e.\ $h=0$) part removed, divided by $u$ to record the
downward move. Collecting,
\[
K(x,u)\,\Phi(x,u) = 1 - \frac{x}{u}\Phi(x,0), \qquad K(x,u):=1-2x-xu-\frac{x}{u}.
\]
The kernel $uK(x,u)=u-2xu-xu^2-x=0$, i.e.\ $xu^2-(1-2x)u+x=0$, has the small root
\[
u_0(x)=\frac{(1-2x)-\sqrt{(1-2x)^2-4x^2}}{2x} = \frac{(1-2x)-\sqrt{1-4x}}{2x},
\]
a power series in $x$ vanishing at $x=0$ ($u_0(x)=x+3x^2+\cdots$). Substituting $u=u_0(x)$ (where
$K=0$, so the left side vanishes) gives $0=1-\frac{x}{u_0(x)}\Phi(x,0)$, i.e.
\[
\Phi(x,0) = \frac{u_0(x)}{x} = \frac{(1-2x)-\sqrt{1-4x}}{2x^2}.
\]
This is the standard kernel-method computation (it is, up to relabelling, the generating function for
Motzkin-type walks with two hold-colors and step-set $\{-1,0,0,+1\}$, ending back at $h=0$); it is
carried out here \emph{as an explicit exercise in the technique on this group's step alphabet}, not
claimed as new combinatorics. It succeeds completely, in the ordinary sense: $\Phi(x,0)$ is algebraic
(a shifted semicircle/Catalan-type series), consistent with the fact that $D$-finiteness (indeed
algebraicity) is exactly what \emph{is} available once there is no value-collapsing equivalence to
impose -- $\Phi$ counts \emph{words}, with no notion of two different profiles representing the same
group element.

\paragraph{Where it stops.} $\Phi(x,0)$ is not $F_T$, or even the one-sided analogue of $F_T$: it counts
words, and (worse) even after passing to \emph{profiles} (grouping words by the pair $(A,K)$ they
realize, which removes the redundancy of choosing which visit deposits which unit, an easy refinement of
the same computation) it still does not identify two profiles $A,A'$ with $A(\zeta)=A'(\zeta)$,
$A\ne A'$. To repair this, one would want a further catalytic variable $v$ carrying the running value
$\sum_{j\le h}a_j\zeta^j$ so that, at the end, one could restrict to $v=$ a fixed target and read off
$F_T$ by extracting a diagonal or coefficient. This is exactly Gap 3's ``carry as catalytic variable.''
Attempting it: extend the step equation to $\Psi(x,u,v)=\sum x^{|w|}u^{h(w)}v^{\text{value}(w)}$, where
an $a^{\pm1}$-step at height $h$ multiplies by $v^{\pm\zeta^h}$. This is \emph{not} a valid catalytic
step for the kernel method as stated: the exponent of $v$ is $\zeta^h$, an irrational (indeed, generically
transcendental-looking, at minimum non-integer) real/complex number depending on $h$, not an integer
increment. The kernel method requires each step to multiply the generating function by $x\cdot(\text{a
fixed Laurent polynomial in the catalytic variables})$, so that the resulting functional equation is
polynomial in $u,v$ with coefficients in $x$; a step that multiplies by $v^{\zeta^h}$ for varying real
$h$-dependent exponent $\zeta^h$ is not expressible this way at all (it is not even a well-defined
formal power series operation in a single variable $v$, since $\zeta^h$ is not an integer). The one
honest repair is to \emph{not} exponentiate $v$ by $\zeta^h$, but instead to keep the profile itself
formal, i.e.\ use $v=t$ (a fresh formal Laurent variable, one exponent per height, matching what $\Phi$
already secretly tracks via $u$): then the ``value'' is only recovered by evaluating the resulting series
at $t=\zeta$, external to the generating-function machinery. This is not a defect of execution; it
identifies the actual shape of the obstruction, stated next.

\section{The obstruction, stated precisely (HEURISTIC diagnosis, not a no-go theorem)}\label{sec:obstruction}

\begin{itemize}
\item The kernel method eliminates an unknown boundary series (here $\Phi(x,0)$, or more generally a
finite tuple of unknown series indexed by finitely many residues/boundary states) by substituting
finitely many small roots $u_i(x)$ of a kernel polynomial $K(x,u)$ of \emph{fixed finite degree in $u$}
(degree $2$ above, matching the two-generator step set $s^{\pm1}$). This works precisely because the
number of unknowns equals the number of independent small roots available.
\item For $F_T$, the analogous ``unknown boundary data'' is not a single series or a finite tuple: it is,
in effect, one unknown quantity \emph{per coset} of $I_T=\ker(\mathrm{ev}_\zeta)\supseteq\mu_T\Z[t^{\pm1}]$
in $\Z[t^{\pm1}]$ (how much of the profile generating function collapses into that coset's target value),
and Lemma~\ref{lem:infinite-carry} shows this coset space is infinite. A degree-$2$ (or any fixed finite
degree $r$) kernel in a catalytic variable $u$ supplies at most $r$ algebraic elimination equations near
$x=0$; it cannot pin down an a-priori-independent family of infinitely many unknown series with only
finitely many equations.
\item This is a cardinality argument, not a full impossibility proof: it would fail to be an obstruction
if the infinitely many coset-unknowns were secretly not independent -- e.g.\ if they satisfied their own
finite-dimensional recursion in the coset index, which is exactly what ``the carry has a rational-base
style bounded automaton'' would have supplied, and which Lemma~\ref{lem:infinite-carry} shows cannot come
from boundedness of the coset space itself (it is infinite), but does not rule out some other
finite-dimensional parametrization of the collapse. No such parametrization is known (this restates
Gap 3), and the present note did not find one after the attempt in \S\ref{sec:kernel}. The diagnosis is
therefore: \emph{the obstruction to a carry-catalytic transfer equation is that the carry ranges over a
provably infinite group (Lemma~\ref{lem:infinite-carry}), and the standard kernel-method elimination step
has a fixed finite algebraic degree, so it cannot see enough of that group to resolve the collapse in one
functional equation of the usual shape} -- this is an argument for why the search of \S\ref{sec:next} of
\texttt{P5\_note.tex} is hard in a specific, structural sense, not a proof that every conceivable
functional-equation route is closed.
\end{itemize}

\section{Weakest points, stated explicitly}\label{sec:weak}

\begin{enumerate}
\item Lemma~\ref{lem:infinite-carry}'s proof of unbounded denominators is correct but was written out by
hand in this session, not checked by an exact-arithmetic script; the inductive claim about
$\beta_n\pmod p$ staying nonzero was checked by the displayed one-line recursion argument only, not by
computing $\beta_n$ symbolically for several $n$. A concrete Rust/sympy check of $\beta_n\bmod p$ for
$(c,e)=(2,1)$, $p=2$, $n\le20$ would upgrade this from ``PROVED by hand, unverified independently'' to
fully checked; not run here (judged not worth the disk/CPU budget for an elementary induction, but flagged
as the honest gap in rigor).
\item \S\ref{sec:kernel}'s successful computation is for the one-sided ($h\ge0$) word-counting series, a
genuinely different (and easier) object than $F_T$ even before the value-collapse issue: it does not even
implement the two-sided TRAVEL cost of Lemma~\ref{lem:travel}, let alone the profile/element distinction.
It demonstrates the technique working on this alphabet, not progress on $F_T$ itself.
\item \S\ref{sec:obstruction}'s cardinality argument is explicitly labelled HEURISTIC and is the weakest
inferential step in this note: ``finitely many algebraic equations cannot determine infinitely many
independent unknowns'' is true when the unknowns are independent, but independence itself is an assumption,
not derived here. A sharper version would need to show the coset-indexed unknown series do not satisfy any
finite recursion in the coset index -- a strictly stronger (and unattempted) claim than
Lemma~\ref{lem:infinite-carry}'s mere infiniteness of the index set.
\item This note does not touch Gap 1 (no known meromorphic continuation of $F_T$ at all) in any way; even
a complete carry-catalytic transfer equation, had one been found, would still need translating into a
statement about analytic continuation and pole count to invoke Lemma~\texttt{lem:reduction} of
\texttt{P5\_note.tex}. That further step was out of scope here and is not attempted.
\item No computation was run for this note beyond hand/mental arithmetic checks (the $\mathrm{TRAVEL}$
sanity check in \S\ref{sec:setup} and the small-root series expansion of $u_0(x)$ in \S\ref{sec:kernel},
both verifiable by hand and not requiring \texttt{wt\_carry} or \texttt{wt\_growth} to be re-run).
\item \textbf{Overall status: Conjecture~\texttt{conj:WT-nonDfinite} remains OPEN.} This note's genuine
new content is Lemma~\ref{lem:infinite-carry} (a PROVED strengthening of Gap 2 from numerics to structure)
and the cardinality diagnosis of \S\ref{sec:obstruction} (a HEURISTIC sharpening of Gap 3 from ``no equation
is known'' to ``a specific, nameable obstruction blocks the standard technique''). Neither closes P5.
\end{enumerate}

\end{document}
